\chapter{常用数据}\label{appendix:常用数据}

    \section{重要常数}\label{app:数据-常数}

        \begin{itemize}
            \item 水的离子积 \(K_\text{W} = \cemh{[H+]} \cdot \cemh{[OH-]}\)：\SI{25}{\degreeCelsius} 时为 \(\num{1E-14}\)，\SI{99}{\degreeCelsius} 时为 \(\num{1E-12}\)。\textcolor{red}{只随温度变化，温度升高增大。}
            \item 常温中性溶液：\(\cemh{[H+]} = \cemh{[OH-]} = \SI{1E-7}{\mole\per\litre}\)，\(\cemh{pH} = 7\)。
            \item \textcolor{blue}{沉淀完全}：溶液中离子浓度 \(< \SI{1E-5}{\mole\per\litre}\)。
            \item \textcolor{blue}{难溶物}：溶解度 \(< \SI{0.01}{g}/100\text{g}\) 水。
        \end{itemize}

    \section{\texorpdfstring{常见弱酸、弱碱的电离常数（\SI{25}{\degreeCelsius}）}{常见弱酸、弱碱的电离常数（25℃）}}\label{app:数据-电离常数}

        数据见正文 \textbf{表~\ref{tab:弱酸-课本3.3}} 及盐类水解一节。

        \par\vspace{0.6em}
        {\footnotesize
        \begin{minipage}[t]{0.47\textwidth}
            \textbf{一元弱酸（或弱碱）}\quad \(K_\text{a}\)（或 \(K_\text{b}\)）
            \par\vspace{0.2em}
            \begin{tabular}{l@{\hspace{1.2em}}r}
                \cemh{HCOOH}（甲酸） & \num{1.8E-4} \\
                \cemh{HNO2}（亚硝酸） & \num{7.2E-4} \\
                \cemh{HF}（氢氟酸） & \num{5.6E-4} \\
                \cemh{CH3COOH}（醋酸） & \num{1.8E-5} \\
                \cemh{HCN}（氢氰酸） & \num{4.0E-10} \\
                \cemh{HClO}（次氯酸） & \num{2.95E-8} \\
                \cemh{NH3 . H2O}（氨水） & \num{1.8E-5}
            \end{tabular}
        \end{minipage}
        \hfill
        \begin{minipage}[t]{0.47\textwidth}
            \textbf{多元弱酸}\quad \(K_\text{a1}\)、\(K_\text{a2}\)（\(K_\text{a3}\)）
            \par\vspace{0.2em}
            \begin{tabular}{l@{\hspace{1.2em}}r}
                \cemh{H2C2O4} & \num{5.6E-2}、\num{5.4E-5} \\
                \cemh{H2SO3} & \num{1.2E-2}、\num{6.2E-8} \\
                \cemh{H3PO4} & \num{7.52E-3} \\
                                & \num{6.23E-8} \\
                                & \num{2.2E-13} \\
                \cemh{H2CO3} & \num{4.2E-7}、\num{4.8E-11} \\
                \cemh{H2S} & \num{5.7E-8}、\num{1.2E-15}
            \end{tabular}
        \end{minipage}}
        \par\vspace{0.6em}

        \textcolor{blue}{酸性强弱：\cemh{H2C2O4} \(>\) \cemh{H2SO3} \(>\) \cemh{HNO2} \(>\) \cemh{HF} \(>\) \cemh{HCOOH} \(>\) \cemh{CH3COOH} \(>\) \cemh{H2CO3} \(>\) \cemh{HCN}。}

        \textcolor{red}{电离常数越小越难电离，其酸根结合 \cemh{H+} 的能力越强：\cemh{CO3^{2-}} \(>\) \cemh{CN-} \(>\) \cemh{CH3COO-} \(>\) \cemh{HSO3-}。}

    \section{\texorpdfstring{常见难溶电解质的溶度积（\SI{25}{\degreeCelsius}）}{常见难溶电解质的溶度积（25℃）}}\label{app:数据-溶度积}

        \par\vspace{0.6em}
        {\footnotesize
        \begin{minipage}[t]{0.47\textwidth}
            \textbf{银盐、硫化物}\quad \(K_\text{sp}\)
            \par\vspace{0.2em}
            \begin{tabular}{l@{\hspace{1.2em}}r}
                \cemh{AgCl} & \num{1.77E-10} \\
                \cemh{AgI} & \num{8.51E-17} \\
                \cemh{Ag2S} & \num{6.69E-50} \\
                \cemh{FeS} & \num{1.59E-19} \\
                \cemh{CuS} & \num{1.27E-36}
            \end{tabular}
        \end{minipage}
        \hfill
        \begin{minipage}[t]{0.47\textwidth}
            \textbf{硫酸盐、碳酸盐、氢氧化物}\quad \(K_\text{sp}\)
            \par\vspace{0.2em}
            \begin{tabular}{l@{\hspace{1.2em}}r}
                \cemh{BaSO4} & \num{1.08E-10} \\
                \cemh{BaCO3} & \num{8.1E-9} \\
                \cemh{CaCO3} & \num{4.96E-9} \\
                \cemh{Fe(OH)3} & \num{2.79E-39} \\
                \cemh{Fe(OH)2} & \num{4.87E-17} \\
                \cemh{Cu(OH)2} & \num{2.2E-20} \\
                \cemh{Mg(OH)2} & \num{8.45E-12} \\
                \cemh{Al(OH)3} & \num{1.1E-33}
            \end{tabular}
        \end{minipage}}
        \par\vspace{0.6em}

        \textcolor{blue}{同类型难溶电解质，\(K_\text{sp}\) 越小溶解度越小、越易生成沉淀；不同类型不能直接比较。}

    \section{常见沉淀的颜色}\label{app:数据-颜色}

        \par\vspace{0.6em}
        {\footnotesize
        \begin{minipage}[t]{0.47\textwidth}
            \textbf{银盐或硫化物}\quad 颜色
            \par\vspace{0.2em}
            \begin{tabular}{l@{\hspace{1.2em}}l}
                \cemh{AgCl} & 白色 \\
                \cemh{AgBr} & 淡黄色 \\
                \cemh{AgI} & 黄色 \\
                \cemh{Ag2S} & 黑色 \\
                \cemh{ZnS} & 白色
            \end{tabular}
        \end{minipage}
        \hfill
        \begin{minipage}[t]{0.47\textwidth}
            \textbf{氢氧化物}\quad 颜色
            \par\vspace{0.2em}
            \begin{tabular}{l@{\hspace{1.2em}}l}
                \cemh{Mg(OH)2} & 白色 \\
                \cemh{Al(OH)3} & 白色 \\
                \cemh{Cu(OH)2} & 蓝色 \\
                \cemh{Fe(OH)3} & 红褐色 \\
                \cemh{Fe(OH)2} & 白色
            \end{tabular}
        \end{minipage}}
        \par\vspace{0.6em}

    \section{常见难溶氢氧化物的沉淀 pH}\label{app:数据-沉淀pH}

        金属离子浓度为 \SI{0.01}{\mole\per\litre}：

        \begin{datas}
            \item \cemh{Fe^{3+}}（\cemh{Fe(OH)3}）\quad 开始沉淀 pH 1.82，完全沉淀 pH 2.82
            \item \cemh{Al^{3+}}（\cemh{Al(OH)3}）\quad 开始沉淀 pH 3.68，完全沉淀 pH 4.68
            \item \cemh{Cu^{2+}}（\cemh{Cu(OH)2}）\quad 开始沉淀 pH 5.17，完全沉淀 pH 6.67
            \item \cemh{Fe^{2+}}（\cemh{Fe(OH)2}）\quad 开始沉淀 pH 6.84，完全沉淀 pH 8.34
            \item \cemh{Mg^{2+}}（\cemh{Mg(OH)2}）\quad 开始沉淀 pH 9.46，完全沉淀 pH 10.96
        \end{datas}

        \textcolor{blue}{用于分离除杂：酸化的 \cemh{CuSO4} 溶液混有 \cemh{Fe^{3+}}，可加 \cemh{CuO} 提高 pH 使 \cemh{Fe^{3+}} 沉淀除去（能与 \cemh{H+} 反应、不引入杂质、过量可过滤）。若还混有 \cemh{Fe^{2+}}，先加 \cemh{H2O2} 氧化。}

    \section{酸碱指示剂的变色范围}\label{app:数据-指示剂}

        \begin{datas}
            \item 甲基橙\quad \numrange{3.1}{4.4}\quad 红 \(\to\) 橙 \(\to\) 黄
            \item 甲基红\quad \numrange{4.4}{6.2}\quad 红 \(\to\) 橙 \(\to\) 黄
            \item 石蕊\quad \numrange{4.5}{8.3}\quad 红 \(\to\) 紫 \(\to\) 蓝
            \item 酚酞\quad \numrange{8.0}{10.0}\quad 无色 \(\to\) 粉红 \(\to\) 红
        \end{datas}

        \textcolor{red}{石蕊变色不明显，滴定中不宜选用；酸滴碱用甲基橙，碱滴酸用酚酞。}

    \section{电解时离子的放电顺序}\label{app:数据-放电顺序}

        \subsubsection{阴极（阳离子得电子）}

            \textcolor{red}{\cemh{Ag+} \(>\) \cemh{Fe^{3+}} \(>\) \cemh{Cu^{2+}} \(>\) \cemh{H+}（酸）\(>\) \cemh{Fe^{2+}} \(>\) \cemh{Zn^{2+}} \(>\) \cemh{H+}（水）\(>\) \cemh{Al^{3+}} \(>\) \cemh{Mg^{2+}} \(>\) \cemh{Na+} \(>\) \cemh{Ca^{2+}} \(>\) \cemh{K+}}

            \textcolor{blue}{注：\cemh{H+} 氢前铝后看浓度。}

        \subsubsection{阳极（阴离子失电子，惰性电极）}

            \textcolor{red}{\cemh{I-} \(>\) \cemh{Br-} \(>\) \cemh{Cl-} \(>\) \cemh{OH-} \(>\) 含氧酸根 \(>\) \cemh{F-}}

            \textcolor{blue}{活性电极（除 \cemh{Pt}、\cemh{Au} 外的金属）优先于阴离子放电——阳极溶解。}

            \textcolor{red}{阳极出氧卤，阴极氢金属。}
